\documentclass[11pt]{article}

\usepackage[margin=0.9in]{geometry}
\usepackage{amsmath,amssymb,amsthm,mathtools,booktabs,enumitem,microtype}
\usepackage{xcolor}
\usepackage[colorlinks=true,linkcolor=blue!50!black,urlcolor=blue!50!black,citecolor=blue!50!black]{hyperref}

\theoremstyle{plain}
\newtheorem{theorem}{Theorem}
\newtheorem{lemma}[theorem]{Lemma}
\newtheorem{proposition}[theorem]{Proposition}
\newtheorem{corollary}[theorem]{Corollary}

\theoremstyle{definition}
\newtheorem{definition}[theorem]{Definition}

\theoremstyle{remark}
\newtheorem{remark}{Remark}

\DeclareMathOperator{\Aut}{Aut}
\DeclareMathOperator{\supp}{supp}
\DeclareMathOperator{\BDC}{BDC}
\newcommand{\eqdist}{\stackrel{\mathcal{D}}{=}}

\title{\bfseries The Symmetry Group of the Binary Deletion Channel}
\author{}
\date{\vspace{-1.6em}%
\small
\begin{tabular}{@{}c@{}}
Target paper: F.~Pernice, \emph{On the Symmetries of the Deletion Channel},\\
\href{https://arxiv.org/abs/2202.03633}{arXiv:2202.03633} (Allerton 2022), Sections~II and~IV, Conjecture~IV.3\\[0.35em]
Companion verification script: \texttt{solutions/bdc\_orbit\_symmetry/explore.py}\\[0.35em]
November 2026
\end{tabular}}

\begin{document}

\maketitle

\begin{abstract}
We prove Conjecture~IV.3 of Pernice (\href{https://arxiv.org/abs/2202.03633}{arXiv:2202.03633}, Allerton~2022) together with the stronger global group conjecture in the surrounding text: for every fixed deletion probability $d \in (0, 1)$, the channel symmetry group $\mathcal{G}_{\BDC_d}$ of the binary deletion channel is precisely the Klein four-group $\{1, C, R, CR\} \cong \mathbb{Z}_2 \times \mathbb{Z}_2$ generated by bitwise complement~$C$ and word reversal~$R$.
Working strictly under Pernice's definition~\texttt{defn:channel-group}---which assumes neither that the output bijection~$h$ preserves length nor that $h = g^{-1}$---we first force $h(A^n) = A^n$ and $h = g^{-1}$ globally from the channel support structure, reducing channel symmetry to preservation of subsequence multiplicities $N_{gx}(gy) = N_x(y)$.
An explicit single-deletion bit-moment formula reconstructs every binary word of length $n \ge 3$ from its single-deletion deck, completing the classification by induction.
Consequently, every orbit in $\{0, 1\}^n$ has size $2$ or $4$ (at most~$2$ for $n \in \{1, 2\}$, attaining~$4$ at $0^{n-1}1$ for every $n \ge 3$).
\end{abstract}

\section{Setup and Reduction to Subsequence Counts}

Let $A := \{0, 1\}$, and let $D = \BDC_d$ denote the binary deletion channel that independently deletes each input bit with fixed probability $d \in (0, 1)$. Following Pernice~\cite[\texttt{defn:channel-group}]{pernice2022} (adapted from Polyanskiy~\cite[Section~VI.A]{polyanskiy2013}), the global channel group of $D$ is
\[
\mathcal{G}_{\BDC_d} := \bigl\{g \in \Aut(A^*) : \exists\, h \in \Aut(A^*),\; h(D(gx)) \eqdist D(x) \text{ and } |g(x)| = |x| \;\; \forall x \in A^*\bigr\},
\]
where $\Aut(A^*)$ is the group of all bijections of $A^*$, and the blocklength-$n$ group $\mathcal{G}_{\BDC_d}^n \subseteq \Aut(A^n)$ is the restriction of $\mathcal{G}_{\BDC_d}$ to $A^n$.

For words $x, y \in A^*$, write $N_x(y)$ for the number of strictly increasing selections of indices in $x$ spelling~$y$ (the subsequence multiplicity). For $|x| = n$ and $|y| = k \le n$, the transition probability is
\[
P_n(d)_{y, x} := \Pr\bigl(D(x) = y\bigr) = N_x(y)\, d^{n - k}(1 - d)^k,
\]
and $P_n(d)_{y, x} = 0$ for $|y| > n$. Since $d \in (0, 1)$, the support $\supp(D(x))$ consists precisely of all subsequences of~$x$, and the unique word of length~$n$ in $\supp(D(x))$ is $x$ itself.

\begin{lemma}[Output bijection forced; subsequence-count equivalence]\label{lem:forced-h}
Fix $d \in (0, 1)$. Let $g \in \Aut(A^*)$ be length-preserving and let $h \in \Aut(A^*)$ satisfy $h(D(gx)) \eqdist D(x)$ for every $x \in A^*$. Then:
\begin{enumerate}[label=\textup{(\arabic*)},nosep]
\item $h(A^n) = A^n$ for every $n \ge 0$, and $h = g^{-1}$ globally on $A^*$.
\item A length-preserving bijection $g \in \Aut(A^*)$ lies in $\mathcal{G}_{\BDC_d}$ if and only if
\begin{equation}\label{eq:subseq}
N_{gx}(gy) = N_x(y) \qquad (\forall x, y \in A^*,\; |y| \le |x|).
\end{equation}
\end{enumerate}
\end{lemma}

\begin{proof}
For every $n \ge 0$, every word $y \in A^{\le n}$ extends to some $x \in A^n$ and arises from~$x$ by deleting the appended suffix; hence $\bigcup_{x \in A^n} \supp(D(x)) = A^{\le n}$. Because $h$ is a bijection with $h(\supp(D(gx))) = \supp(D(x))$ and $g(A^n) = A^n$, taking unions over $x \in A^n$ gives
\[
h(A^{\le n}) = h\Bigl(\bigcup_{x \in A^n} \supp(D(gx))\Bigr) = \bigcup_{x \in A^n} \supp(D(x)) = A^{\le n} \qquad (\forall n \ge 0).
\]
Taking successive set differences ($A^0 = A^{\le 0}$ and $A^n = A^{\le n} \setminus A^{\le n-1}$) forces $h(A^n) = A^n$ for all $n \ge 0$. For any $x \in A^n$, intersecting $h(\supp(D(gx))) = \supp(D(x))$ with $A^n$ yields
\[
\{h(gx)\} = h\bigl(\supp(D(gx)) \cap A^n\bigr) = \supp(D(x)) \cap A^n = \{x\},
\]
so $h(gx) = x$ for every $x \in A^*$, i.e., $h = g^{-1}$.

Substituting $h = g^{-1}$ into $h(D(gx)) \eqdist D(x)$ gives $P_n(d)_{gy, gx} = P_n(d)_{y, x}$ for all $|y| = k \le n = |x|$. Since $d^{n-k}(1 - d)^k > 0$ depends only on $(n, k)$ and $g$ preserves length, cancelling this positive factor yields~\eqref{eq:subseq}. Conversely,~\eqref{eq:subseq} implies $P_n(d)_{gy, gx} = P_n(d)_{y, x}$ and hence $g^{-1}(D(gx)) \eqdist D(x)$ at every fixed $d \in (0, 1)$, including $d = 1/2$.
\end{proof}

\section{Single-Deletion Reconstruction and Classification}

\begin{lemma}[Single-deletion moment reconstruction]\label{lem:reconstruct}
For any binary word $x = x_1 \cdots x_n \in \{0, 1\}^n$ (viewing bits as integers in $\{0, 1\}$) and every $1 \le i \le n - 1$,
\begin{equation}\label{eq:moments}
s_i(x) := \sum_{y \in \{0, 1\}^{n-1}} N_x(y)\, y_i = (n - i) x_i + i\, x_{i+1}.
\end{equation}
Consequently, for every $n \ge 3$, the single-deletion multiplicity deck $\{N_x(y)\}_{y \in \{0, 1\}^{n-1}}$ uniquely determines~$x$.
\end{lemma}

\begin{proof}
Consider the $n$ single-position deletions of $x$ with multiplicity. The $n - i$ deletions strictly after position~$i$ leave $x_i$ at position~$i$, while the $i$ deletions at or before position~$i$ shift $x_{i+1}$ into position~$i$. Summing the $i$th bit across all $n$ outcomes yields~\eqref{eq:moments}.
For $n \ge 3$, we have $n - 1 \ge 2$ and $x_2 \in \{0, 1\}$, so $s_1(x) = (n - 1)x_1 + x_2 \ge n - 1$ if and only if $x_1 = 1$. Once $x_i$ is recovered, $x_{i+1} = (s_i(x) - (n - i)x_i)/i$ recovers $x$ inductively. (At $n = 2$, $01$ and $10$ share the same single-deletion deck.)
\end{proof}

(For the broader combinatorial theory of reconstructing words from subsequence decks, see Manvel, Meyerowitz, Schwenk, Smith, and Stockmeyer~\cite{manvel1991}.)

\begin{theorem}[Classification of $\mathcal{G}_{\BDC_d}$ and orbit sizes]\label{thm:main}
Fix any deletion probability $d \in (0, 1)$. Let $C$ denote bitwise complement and $R$ denote word reversal on $A^*$. Then:
\begin{enumerate}[label=\textup{(\arabic*)},nosep]
\item $\mathcal{G}_{\BDC_d} = \{1, C, R, CR\} \cong \mathbb{Z}_2 \times \mathbb{Z}_2$.
\item The blocklength-$n$ restriction $\mathcal{G}_{\BDC_d}^n$ has order~$2$ for $n = 1$ and order~$4$ for every $n \ge 2$.
\item For every $n \ge 1$ and $x \in \{0, 1\}^n$, the orbit of~$x$ is $\{x, Cx, Rx, CRx\}$ (deduplicated), with cardinality~$2$ if $Rx \in \{x, Cx\}$ and cardinality~$4$ otherwise. In particular, every orbit has size at most~$4$ (maximally~$2$ for $n \in \{1, 2\}$, and $4$ for every $n \ge 3$, attained at $0^{n-1}1$).
\end{enumerate}
\end{theorem}

\begin{proof}
Both $C$ and $R$ are commuting length-preserving involutions satisfying~\eqref{eq:subseq}: complement preserves the exact selected index sets, and reversal reflects selected index sets and reads them in reverse order. Thus $\{1, C, R, CR\} \subseteq \mathcal{G}_{\BDC_d}$.

Conversely, let $g \in \mathcal{G}_{\BDC_d}$. Its restriction to $A^1 = \{0, 1\}$ is either the identity or the swap; composing with $C$ if needed yields a symmetry fixing $0$ and $1$. By~\eqref{eq:subseq} with $|y| = 1$, this symmetry preserves Hamming weight $N_x(1)$, so on $A^2$ it fixes $00$ and $11$ and permutes $\{01, 10\}$. Composing with $R$ if needed produces a symmetry $f \in \mathcal{G}_{\BDC_d}$ that fixes $A^0 \cup A^1 \cup A^2$ pointwise.

We show by induction on $n \ge 2$ that $f$ fixes $A^n$ pointwise. The case $n = 2$ holds by construction. Assuming $n \ge 3$ and $f|_{A^{n-1}} = \mathrm{id}$, equation~\eqref{eq:subseq} gives, for every $x \in A^n$ and every $y \in A^{n-1}$,
\[
N_{fx}(y) = N_{fx}(fy) = N_x(y).
\]
By Lemma~\ref{lem:reconstruct}, $fx = x$ for every $x \in A^n$. Induction yields $f = \mathrm{id}$ globally on $A^*$, whence $g \in \{1, C, R, CR\}$.

At $n = 1$, $R$ acts trivially while $C$ swaps $0$ and $1$, so $|\mathcal{G}_{\BDC_d}^1| = 2$ and there is one orbit of size~$2$. For $n \ge 2$, the constant words $0^n, 1^n$ distinguish $\{C, CR\}$ from $\{1, R\}$, and $0^{n-1}1$ distinguishes $R$ from~$1$, giving $|\mathcal{G}_{\BDC_d}^n| = 4$. Because $Cx \neq x$ for every nonempty $x$, every orbit has size $2$ or $4$, with size $2$ iff $Rx \in \{x, Cx\}$ (always true for $n = 2$). For every $n \ge 3$, $x = 0^{n-1}1$ satisfies $Rx = 10^{n-1} \notin \{0^{n-1}1, 1^{n-1}0\}$, attaining orbit size~$4$.
\end{proof}

\begin{remark}[Subword-order automorphism route]
Once Lemma~\ref{lem:forced-h} establishes $h = g^{-1}$, positive-probability transitions encode the subsequence partial order $y \preceq x$, so $g$ is an order automorphism of $(A^*, \preceq)$. Lemma~3.8 of Halfon, Schnoebelen, and Zetzsche~\cite{halfon2017} classifies all automorphisms of $(A^*, \preceq)$ as alphabet permutations optionally composed with reversal, giving an independent second half of Theorem~\ref{thm:main} for $|A| = 2$.
\end{remark}

\begin{remark}[Endpoint cases and global vs.\ single-blocklength scope]
At $d = 0$ ($D(x) = x$, with $h = g^{-1}$) and at $d = 1$ ($D(x) = \varepsilon$, with $h = \mathrm{id}$), every length-preserving bijection belongs to $\mathcal{G}_{\BDC_d}$, so $\mathcal{G}_{\BDC_d}^n \cong S_{2^n}$ and the single orbit has size $2^n$; thus $d \in (0, 1)$ is essential.
Note also that Pernice's~\texttt{defn:channel-group} is global on $A^*$, which is what Theorem~\ref{thm:main} classifies. Even for a single fixed blocklength~$n$, every $y \in A^k$ ($0 \le k \le n$) appears in $\supp(D(x))$ for exactly $\sum_{j=k}^n \binom{n}{j}$ inputs $x \in A^n$---a count strictly decreasing in~$k$---so any output bijection matching the output distributions of $A^n$ still satisfies $h(A^k) = A^k$ and $h|_{A^n} = g^{-1}$; classifying standalone single-length channel groups at fixed~$n$ is a separate question beyond Pernice's conjecture.
\end{remark}

\section{Exact Finite Verification}

The standard-library script \texttt{solutions/bdc\_orbit\_symmetry/explore.py} computes every transition column in exact factored form $N_x(y)\,d^{n - |y|}(1 - d)^{|y|}$ across all 8{,}190 nonempty binary words of lengths $1 \le n \le 12$. It checks layer binomial sums, single-deletion run-length multiplicities, the moment identity~\eqref{eq:moments}, exact moment reconstruction for every $n \ge 3$, and exact equality between unlabelled layer-multiplicity histograms and $\{1, C, R, CR\}$ orbits (Table~\ref{tab:orbits}). All assertions pass with exit code~$0$.

\begin{table}[h]
\centering
\begin{tabular}{@{}rrrrr@{\hspace{2.2em}}rrrrr@{}}
\toprule
$n$ & $2^n$ & Size-2 & Size-4 & Max & $n$ & $2^n$ & Size-2 & Size-4 & Max \\
\midrule
1 &  2 &  1 &  0 & 2 &  7 &  128 &  8 &  28 & 4 \\
2 &  4 &  2 &  0 & 2 &  8 &  256 & 16 &  56 & 4 \\
3 &  8 &  2 &  1 & 4 &  9 &  512 & 16 & 120 & 4 \\
4 & 16 &  4 &  2 & 4 & 10 & 1024 & 32 & 240 & 4 \\
5 & 32 &  4 &  6 & 4 & 11 & 2048 & 32 & 496 & 4 \\
6 & 64 &  8 & 12 & 4 & 12 & 4096 & 64 & 992 & 4 \\
\bottomrule
\end{tabular}
\caption{Exact orbit counts of $\{0, 1\}^n$ under $\mathcal{G}_{\BDC_d}$ verified by \texttt{explore.py}.}\label{tab:orbits}
\end{table}

\begin{thebibliography}{9}

\bibitem{halfon2017}
S.~Halfon, P.~Schnoebelen, and G.~Zetzsche.
Decidability, complexity, and expressiveness of first-order logic over the subword ordering.
In \emph{Proceedings of the 32nd Annual ACM/IEEE Symposium on Logic in Computer Science (LICS 2017)},
pp.~1--12, \href{https://arxiv.org/abs/1701.07470}{arXiv:1701.07470}, 2017.

\bibitem{manvel1991}
B.~Manvel, A.~Meyerowitz, A.~J.~Schwenk, K.~W.~Smith, and P.~K.~Stockmeyer.
Reconstruction of sequences.
\emph{Discrete Mathematics}, 94(3):209--219, 1991.

\bibitem{pernice2022}
F.~Pernice.
On the symmetries of the deletion channel.
In \emph{58th Annual Allerton Conference on Communication, Control, and Computing},
\href{https://arxiv.org/abs/2202.03633}{arXiv:2202.03633}, 2022.

\bibitem{polyanskiy2013}
Y.~Polyanskiy.
Saddle point in the minimax converse for channel coding.
\emph{IEEE Transactions on Information Theory}, 59(5):2576--2595, 2013.

\end{thebibliography}

\end{document}
